% Generado: no editar cifras manualmente.
\begin{table}[!htbp]
\centering
\caption{Especificidad del bloque de donativos públicos por ronda}\label{tab:catalogo-instrumentos}
\small\setlength{\tabcolsep}{4pt}
\renewcommand{\arraystretch}{1.15}
\begin{tabular}{lp{3.3cm}p{8cm}}
\toprule
\textbf{Año} & \shortstack[c]{\textbf{Recepción}\\\textbf{A propia}} & \shortstack[c]{\textbf{Identificación}\\\textbf{de programa}} \\
\midrule
2001 & k76a & sin programa nombrado \\
2003 & k79a & sin programa nombrado \\
2012 & k79a\_12 & sin programa nombrado \\
2015 & k79a\_15 & sin programa nombrado \\
2018 & k79a\_18 & 1: PROAGRO Productivo; 2: PROSPERA; 3: INAPAM; 4: Otra institución: tipo no resuelto \\
2021 & k79a\_21 & 1: Producción para el Bienestar; 2: Becas para el Bienestar Benito Juárez; 3: INAPAM; 4: Otra institución: tipo no resuelto; 5: Programa para el Bienestar de las Personas Adultas Mayores \\
\bottomrule
\end{tabular}
\par\medskip
\parbox{\linewidth}{\footnotesize Fuente: contratos documentados de cuestionarios y libros ENASEM; no se recalculan respuestas. Propio refiere a la rama económica del beneficiario validada, no necesariamente al informante. Cónyuge, códigos, slots y montos se conservan en la matriz completa. Antes de2018 no se observa nombre de programa en este bloque. La categoría otro y las fuentes públicas K permanecen indeterminadas.}
\end{table}
